\chapter{Construction of It\^o's Integral with Respect to Continuous \texorpdfstring{$L^2$}{L2}-Bounded Martingales}

Let \(S\leq T\) be stopping times and let \(Z\) be bounded and
\(\F_S\)-measurable. Consider the elementary predictable process
\[
    H_s=Z\mathbf 1_{(S,T]}(s).
\]
For a continuous, \(L^2\)-bounded martingale \(M\), define
\[
    \int_0^tH_s\,dM_s
    :=Z\bigl(M_{t\wedge T}-M_{t\wedge S}\bigr).
\]
We also write this stochastic integral as \((H\cdot M)_t\).

\begin{proposition}
The process
\[
    (H\cdot M)_t=Z\bigl(M_{t\wedge T}-M_{t\wedge S}\bigr)
\]
is an adapted, \(L^2\)-bounded martingale.
\end{proposition}

\begin{proof}
The process is adapted: before \(S\) it is zero, while for \(t\geq S\),
\(Z\) is \(\F_S\subseteq\F_t\)-measurable.

For \(L^2\)-boundedness, conditional Jensen and optional sampling give
\[
    \|M_{t\wedge T}\|_2\leq\|M_t\|_2,
    \qquad
    \|M_{t\wedge S}\|_2\leq\|M_t\|_2.
\]
Thus
\[
    \|(H\cdot M)_t\|_2
    \leq2\|Z\|_\infty\|M_t\|_2,
\]
so \(H\cdot M\) is \(L^2\)-bounded.

It remains to prove the martingale property. Let \(R\) be an arbitrary
bounded stopping time. For \(A\in\F_S\), define
\[
    S_A:=S\mathbf 1_A+\infty\mathbf 1_{A^c},
    \qquad
    T_A:=T\mathbf 1_A+\infty\mathbf 1_{A^c}.
\]
Both are stopping times, and optional sampling applied to the bounded
stopping times \(S_A\wedge R\leq T_A\wedge R\) gives
\[
    \E[M_{T_A\wedge R}]=\E[M_{S_A\wedge R}].
\]
Outside \(A\), both stopped variables equal \(M_R\); hence
\[
    \E\bigl[(M_{T\wedge R}-M_{S\wedge R});A\bigr]=0.
\]
Since this holds for every \(A\in\F_S\),
\[
    \E[M_{T\wedge R}-M_{S\wedge R}\mid\F_S]=0.
\]
Multiplying by the bounded \(\F_S\)-measurable variable \(Z\), then taking
expectations, gives \(\E[(H\cdot M)_R]=0=(H\cdot M)_0\). The converse
optional-sampling criterion now implies that \(H\cdot M\) is a martingale.
\end{proof}
\section{Lecture 17}

For more general $H$, we can take linear combinations of the earlier form:
\[
H = \sum_{i=1}^n Z_{i-1} \mathbf{1}_{(T_{i-1}, T_i]},
\]
where $0 \le T_0 < \cdots < T_n$ are stopping times. We will then define 
\[
(H \cdot M)_t = \sum Z_{i-1}(M_{T_i \wedge t} - M_{T_{i-1} \wedge t}).
\]

\begin{center}
%html-figure: ito-construction-roadmap.html
\begin{tikzpicture}[
    >=Stealth,
    box/.style={
        draw,
        rounded corners=8pt,
        minimum width=5.4cm,
        minimum height=2.6cm,
        align=center,
        inner sep=8pt
    },
    roadmaptext/.style={
        font=\small\itshape,
        align=center
    }
]

% main objects
\node[box] (integrands) at (0,0) {
    \textbf{Elementary integrands}\\[2mm]
    $\displaystyle
    H=\sum_{i=1}^n Z_{i-1}\mathbf{1}_{(T_{i-1},T_i]}
    $
};

\node[box] (outputs) at (8.2,0) {
    \textbf{Resulting stochastic processes}\\[2mm]
    $\displaystyle
    H\cdot M
    =
    \left\{
        \int_0^t H_s\,dM_s
    \right\}_{t\geq 0}
    $
};


% steps 1 and 2
\node[
    roadmaptext,
    text width=5.4cm
] at (0,-3.15) {
    \textbf{1.} Give the integrands a norm\\
    and a measurable structure.
};

\node[
    roadmaptext,
    text width=5.4cm
] at (8.2,-3.15) {
    \textbf{2.} Identify $L^2$-bounded martingales\\
    with their terminal values in $L^2(\F_\infty)$.
};

% step 3
\node[
    roadmaptext,
    text width=11.5cm
] (density) at (4.1,-5.15) {
    \textbf{3.} Show that the elementary integrands are dense\\
    in the completed integrand space $L^2(M)$.
};

% step 4
\draw[->, thick] (4.1,-5.95) -- (4.1,-6.75);

\node[
    roadmaptext,
    text width=11.5cm
] at (4.1,-7.25) {
    \textbf{4.} Prove the It\^o isometry, then extend the integral\\
    uniquely to every predictable process in $L^2(M)$.
};

\end{tikzpicture}
\end{center}



For the purpose of (2), 
we will use the fact that an $L^2$-bounded martingale can be identified with elements in $L^2(\F_\infty)$. 
Recall that by the martingale convergence theorem, if $\{M_t\}$ is $L^2$-bounded, then $M_t \overset{L^2}{\to} M_\infty \in L^2(\F_\infty)$. Then, $M_t = \E(M_\infty \mid \F_t)$. 

Thus, we have a one-to-one relation 
\[
\{M_t\}_{t \ge 0} \longleftrightarrow M_\infty \in L^2(\F_\infty). 
\]

For (1), note that
\begin{align*}
\E((H \cdot M)_\infty^2) &= \E\left(\left(\sum Z_{i-1}(M_{T_i} - M_{T_{i-1}})\right)^2\right) \\ 
&= \E \left( \sum Z_{i-1}^2 (M_{T_i}-M_{T_{i-1}})^2\right)
\end{align*}
Since $M_t^2 - [M]_t$ is a uniformly integrable martingale, for stopping times $S \le T$, we get 
\begin{align*}
\E((M_T-M_S)^2 \mid \F_S) &= \E[M_T^2-M_S^2 \mid \F_S] \\ 
&= \E([M]_T - [M]_S \mid \F_S)
\end{align*}
Thus, 
\begin{align*}
\E((H \cdot M)_\infty^2) &= \E\left(\sum Z_{i-1}^2([M]_{T_i} - [M]_{T_{i-1}})\right)  \\ 
&= \E \left(\int_0^\infty H_s^2 d[M]_s \right)
\end{align*}
We will then define 
\[
\|H\|_M  = \E \left(\int_0^\infty H_s^2 d[M]_s \right)^{1/2}
\]

For (3), let $b\mcal{E}$ be the set of all processes $H$ with the form 
\[
H = \sum_{i=1}^n Z_{i-1} \mathbf{1}_{(T_{i-1}, T_{i}]}, \quad Z_i \in \F_{T_i}
\]
the b stands for bounded, the E for elementary. 

\begin{definition}
The $\sigma$-algebra generated by $b\mcal{E}$ is 
\[
\mcal{P} = \sigma(b \mcal{E}),
\]
which is called the predictable/previsible $\sigma$-algebra.
\end{definition}

A process is called predictable/previsible if it is a $\mcal{P}$-measurable mapping from $[0, \infty) \times \Omega \to (\R, \mcal{B}(\R))$.

What processes are predictable? By monotone class theorem, all the Càglàd (continuous from the left with right limits existing) functions (in our case, Stochastic processes) are.

Define 
\[
L^2(M) = \left\{\text{predictable processes } H \text{ such that } \|H\|_M < \infty \right\}.
\]
We have defined $H \cdot M$ for $H \in b\mcal{E}$. Our goal is to show that $\overline{b \mcal{E}} = L^2(M)$, where the closure is taken under the earlier norm. 
The proof is omitted. (The idea is to use a version of the monotone class theorem.)

For (4), we have seen that 
\[
\E((H \cdot M)_\infty^2) =
 \E \left(\int_0^\infty H_s^2 d[M]_s \right)
\]
One side is $\|(H \cdot M)_\infty\|_2^2$, and the other is $\|H\|_M^2$. This means that the mapping $I \colon b\mcal{E} \cap L^2(M) \to L^2(\F_\infty)$ defined by $I(H) = (H \cdot M)_\infty$ is an isometry. Since $b\mcal{E}$ is dense in $L^2(M)$, this isometry extends uniquely to the closure $\overline{b \mcal{E}} = L^2(M)$.
